% Block 5 of the slides (05_structure.tex).
\section{Project structure, tests and continuous integration}\label{sec:structure}

\onslides{which code made Figure~3; the DrWatson layout; the code as one part of the
research project; \texttt{src} computes, scripts orchestrate; every result knows its
commit; the environment is part of the result; tests; continuous integration; the same
principles for MATLAB, instruments and FPGA.}

The goal on the first slide is the test of this block: \emph{for every figure, you can say
which code, data and parameters made it}. An agent makes this harder, because it creates
files faster than anyone names them well, and easier, because it follows a fixed layout
reliably once the layout exists.

\subsection{The DrWatson layout}

DrWatson \rf{drwatson,drwatson-docs} is a Julia package that fixes a project layout and
provides a few functions for naming, saving and loading results.
\cmd{initialize\_project("DrivenSpin")} creates the folders and a git repository:
\begin{lstlisting}
DrivenSpin/
|-- src/            # the physics: functions, no hard-coded parameters
|-- scripts/        # short runs that call src/ and save results
|-- test/           # runtests.jl: the physics checks
|-- data/
|   |-- exp_raw/    # measurements: read-only, never modified
|   |-- exp_pro/    # processed data
|   `-- sims/       # simulation results (generated)
|-- plots/          # figures (generated)
|-- Project.toml    # dependencies
|-- Manifest.toml   # exact versions
`-- AGENTS.md       # instructions for the agent
\end{lstlisting}
Functions such as \cmd{datadir("sims", "rabi")} and \cmd{plotsdir()} return absolute
paths into these folders, so a script works from any working directory. A project with
this layout reads the same to you, your colleagues and the agent. The raw-data folder is
read-only by convention; \cmd{AGENTS.md} repeats the rule.

\textbf{The code is one part of the project.} A research project also has a manuscript,
notes and talks, with other collaborators and another history. The code gets its own
repository inside the project folder, so it can be shared, tested and published on its
own; in the course, \cmd{DrivenSpin/} always means that code repository.

\subsection{Scripts and \texttt{produce\_or\_load}}

The slide's script (it becomes \cmd{scripts/rabi.jl} in \cmd{DrivenSpin/}):
\begin{lstlisting}
using DrWatson
@quickactivate "DrivenSpin"          # this project's environment, from any folder
using DrivenSpin                     # the module in src/

for c in dict_list(Dict(:Δ => [0.0, 2π], :Ω => 2π))   # rad/μs; one run per Δ
    produce_or_load(c, datadir("sims", "rabi")) do c
        H = hamiltonian(c[:Δ], c[:Ω])
        t = range(0, 3, length = 301)                  # μs
        Dict("t" => t, "Pe" => excited_population.(Ref(H), t))
    end
end
\end{lstlisting}
\begin{placeholder}[Placeholder: \texttt{DrivenSpin/scripts/rabi.jl}]
Replace the listing above by the real file once D3 exists (it should be identical; the
slide was written first).
\end{placeholder}

Line by line:
\begin{itemize}
  \item \cmd{@quickactivate "DrivenSpin"} activates the project environment (the
        \cmd{Project.toml} of \cmd{DrivenSpin/}) wherever the script is started from.
  \item \cmd{dict\_list} expands a dictionary whose values may be vectors into all
        combinations: here two configurations, $\Delta = 0$ and $\Delta = 2\pi$, both with
        $\Omega = 2\pi$.
  \item \cmd{produce\_or\_load(f, config, path)} builds a file name from \cmd{config} with
        \cmd{savename}, and returns \cmd{(data, file)}. If the file exists, it is loaded;
        otherwise \cmd{f(config)} runs and its result is saved \rf{drwatson-docs}. The
        \cmd{do} block is Julia syntax for passing \cmd{f} as an anonymous function.
  \item \cmd{savename} sorts the keys, joins \cmd{key=value} pairs with \cmd{\_} and rounds
        numbers to three significant digits by default, which gives
        \cmd{Δ=0.0\_Ω=6.28.jld2}.
  \item For \cmd{.jld2} files the result is saved with \cmd{tagsave}, which adds the key
        \cmd{gitcommit}: the output of \cmd{git describe} for the project repository
        \rf{drwatson-docs}.
\end{itemize}
If the repository had uncommitted changes, \cmd{gitcommit} ends in \cmd{-dirty}: the
code that made the file is not in any commit, so that version is lost. DrWatson can also
store the uncommitted diff under the key \cmd{gitpatch} (keyword \cmd{storepatch}, off by
default). Hence the rule on the slide: \emph{commit first, then run}.

Two cautions for questions. First, \cmd{produce\_or\_load} decides by file name only: if
you change the code and rerun, it loads the old result. Delete the file or pass
\cmd{force = true}. Second, three significant digits mean $\Delta = 6.283$ and
$\Delta = 6.284$ share a file name; choose parameters, or \cmd{savename} options, so that
names stay distinct.

\subsection{The environment is part of the result}

\cmd{Project.toml} lists what you depend on and is written by \cmd{Pkg.add}.
\cmd{Manifest.toml} records the exact version of every package, including indirect
dependencies, and is never edited by hand. With both committed,
\cmd{Pkg.instantiate()} reproduces your environment on another machine. An agent that
``fixes'' an error with \cmd{Pkg.update()} changes package versions, and numbers may change
with them; a \cmd{Manifest.toml} in a diff you did not expect needs an explanation.

\subsection{Tests}

\begin{lstlisting}
using DrivenSpin, Test

@testset "Rabi" begin
    Δ, Ω = 0.3, 2π                                  # rad/μs
    ΩR = sqrt(Ω^2 + Δ^2)
    H = hamiltonian(Δ, Ω)
    for t in range(0, 3, length = 31)               # μs
        @test excited_population(H, t) ≈ Ω^2 / ΩR^2 * sin(ΩR * t / 2)^2 atol = 1e-10
    end
    @test excited_population(hamiltonian(0.0, Ω), π / Ω) ≈ 1    # π-pulse
end
\end{lstlisting}
\begin{placeholder}[Placeholder: \texttt{DrivenSpin/test/runtests.jl}]
Replace by the stage-1 tests of D3 once they exist.
\end{placeholder}

\cmd{@testset} groups tests under a name that says which physical claim they check;
\cmd{@test} records a pass or a failure without stopping \rf{julia-test}. The operator
\lstinline|≈| is \cmd{isapprox}, which by default compares with a \emph{relative} tolerance
of $\sqrt{\epsilon_{\mathrm{mach}}} \approx 1.5 \times 10^{-8}$, where $\epsilon_{\mathrm{mach}}$ is the
double-precision machine epsilon, and no absolute tolerance. Near zero a relative tolerance is
meaningless: \verb|isapprox(1e-12, 0.0)| is \cmd{false}. Hence the \cmd{atol} on the
slide wherever the expected value can vanish ($\Pe(0) = 0$). Run the tests with
\verb|julia --project -e 'using Pkg; Pkg.test()'|, before every commit; keep them under a
minute.

A good test checks \emph{physics}, not the implementation (a formula, a limit, a
conservation law); it is deterministic (fixed random seed); it is fast; and it fails when
the physics is wrong. The last property is the one to check deliberately: break the code
and watch the test fail. A test that cannot fail checks nothing.

\subsection{Continuous integration}

Continuous integration (CI) runs the tests on a clean machine at every push. GitHub
Actions reads workflow files in \cmd{.github/workflows/}. The slide's file:
\begin{lstlisting}
name: CI
on: [push, pull_request]
jobs:
  test:
    runs-on: ubuntu-latest
    steps:
      - uses: actions/checkout@v4             # get the repository
      - uses: julia-actions/setup-julia@v2    # install Julia
      - uses: julia-actions/cache@v2          # cache packages between runs
      - uses: julia-actions/julia-buildpkg@v1 # instantiate the environment
      - uses: julia-actions/julia-runtest@v1  # Pkg.test()
\end{lstlisting}
Each \cmd{uses:} line runs a published action \rf{julia-actions}. The result appears as a
check mark or a cross next to the commit, under the \emph{Actions} tab. CI catches
``works on my laptop'' (a missing dependency, a file that was never committed) and makes
sure nobody, including the agent, forgets to run the tests. The file is identical for any
Julia project. The workflow runs only when \cmd{.github/} sits at the root of a
repository, which is why \cmd{DrivenSpin/} carries its own copy.

\subsection{Your tools, same principles}

\begin{center}
\begin{tabularx}{\linewidth}{@{}l X@{}}
\toprule
\textbf{Principle} & \textbf{In the lab} \\
\midrule
Version control & Analysis scripts, instrument settings, HDL sources of FPGA designs. \\
Tests & MATLAB's \cmd{matlab.unittest} framework \rf{matlab-unittest}; simulation
  testbenches for FPGA designs. \\
Raw data is read-only & The agent never writes where measurements land. \\
Traceability & Save the commit hash and the settings with every measurement run. \\
\bottomrule
\end{tabularx}
\end{center}

\begin{qa}
\Q{Do I have to use DrWatson?}
No. The principles are a fixed layout, results tied to a commit, and a pinned
environment. DrWatson gives them to Julia projects in a few lines.

\Q{Why not commit the results too?}
They can be regenerated from code, data and parameters, which are committed; large
binary files bloat the history. Archive final results that a paper depends on in a data
repository with a DOI.

\Q{How do I test code with random numbers?}
Fix the seed (\cmd{Random.seed!}) in the test, and test statistical claims with a
tolerance derived from the expected spread (Section~\ref{sec:verification}).

\Q{Does CI cost money?}
GitHub Actions is free for public repositories; private repositories have a monthly
quota of free minutes.
\end{qa}

\textbf{Read next:} \emph{Good enough practices in scientific computing} \rf{wilson-good};
the DrWatson documentation \rf{drwatson-docs}, ``Workflow tutorial''; the Scientific
Python Development Guide \rf{scipy-dev} for the same ideas in another language.
